% =====================================================================
%  negocio.tex  --  Documento de Negocio de LEIP
%  ARCHIVO UNICO: no necesita ningun otro archivo.
%  Compilar en Overleaf con pdfLaTeX.
% =====================================================================
\documentclass[11pt,a4paper]{article}

% ---------------------------------------------------------------------
%  DATOS DEL PROYECTO  (tocá solo esto para actualizar el documento)
% ---------------------------------------------------------------------
\newcommand{\proyectoNombre}{LEIP}
\newcommand{\proyectoNombreLargo}{LatAm Economic Intelligence Platform}
\newcommand{\proyectoUniversidad}{Universidad Cat\'olica del Uruguay}
\newcommand{\proyectoFacultad}{Facultad de Ingenier\'ia y Tecnolog\'ias}
\newcommand{\proyectoCarrera}{Ingenier\'ia en Inform\'atica}
\newcommand{\docVersion}{0.2}
\newcommand{\docTitulo}{Documento de Negocio}
\newcommand{\docFecha}{Septiembre 2026}
\newcommand{\equipoLista}{%
  Sebastian Forische \and
  Armando Hernandez de la Vega \and
  Luana Acu\~na \and
  Bruno Sebastian Olivera Viera Gomez%
}
\newcommand{\modMacro}{Macro Expectations Intelligence}
\newcommand{\modPrecios}{Prices \& Inflation Intelligence}
\newcommand{\modActividad}{Real Activity Intelligence}

% ---------------------------------------------------------------------
%  PREAMBULO  (estilos y comandos; normalmente no hace falta tocarlo)
% ---------------------------------------------------------------------
% =====================================================================
%  preamble.tex — Paquetes, estilo y comandos. Compartido por TODOS los docs.
% =====================================================================

\usepackage[utf8]{inputenc}
\usepackage[T1]{fontenc}
% microtype requiere fuentes escalables; va junto con lmodern.
\IfFileExists{lmodern.sty}{\usepackage{lmodern}\usepackage{microtype}}{}

% --- Idioma -------------------------------------------------------------
% Si babel-spanish está instalado (Overleaf, TeX Live completo) lo usa.
% Si no, cae a inglés y renombra los títulos a mano para no romper el build.
\IfFileExists{spanish.ldf}{%
  \usepackage[spanish,es-tabla,es-noquoting]{babel}%
}{%
  \usepackage[english]{babel}%
  \renewcommand{\contentsname}{\'Indice}%
  \renewcommand{\listfigurename}{\'Indice de figuras}%
  \renewcommand{\listtablename}{\'Indice de tablas}%
  \renewcommand{\tablename}{Tabla}%
  \renewcommand{\figurename}{Figura}%
  \renewcommand{\refname}{Referencias}%
  \renewcommand{\abstractname}{Resumen}%
}

% --- Página -------------------------------------------------------------
\usepackage[a4paper,margin=2.5cm,headheight=15pt]{geometry}
\usepackage{fancyhdr}
\usepackage{parskip}

% --- Tablas y listas ----------------------------------------------------
\usepackage{booktabs}
\usepackage{tabularx}
\usepackage{longtable}
\usepackage{array}
\usepackage{enumitem}
\setlist{noitemsep,topsep=4pt}

% --- Gráficos y color ---------------------------------------------------
\usepackage{graphicx}
\usepackage{xcolor}
\usepackage{tcolorbox}
\tcbuselibrary{breakable}

% --- Enlaces ------------------------------------------------------------
\usepackage[hidelinks]{hyperref}
\usepackage{url}
\urlstyle{same}

% --- Paleta -------------------------------------------------------------
\definecolor{leipAzul}{HTML}{1F4E79}
\definecolor{leipGris}{HTML}{5A5A5A}
\definecolor{leipAmbar}{HTML}{B8860B}
\definecolor{leipAmbarBg}{HTML}{FFF8E1}

% --- Encabezado y pie ---------------------------------------------------
\pagestyle{fancy}
\fancyhf{}
\fancyhead[L]{\small\color{leipGris}\proyectoNombre{} --- \docTitulo}
\fancyhead[R]{\small\color{leipGris} v\docVersion}
\fancyfoot[C]{\small\thepage}
\renewcommand{\headrulewidth}{0.4pt}

% --- Títulos de sección -------------------------------------------------
\usepackage{titlesec}
\titleformat{\section}{\Large\bfseries\color{leipAzul}}{\thesection}{0.6em}{}
\titleformat{\subsection}{\large\bfseries\color{leipAzul}}{\thesubsection}{0.6em}{}
\titleformat{\subsubsection}{\normalsize\bfseries\color{leipGris}}{\thesubsubsection}{0.6em}{}

% =====================================================================
%  COMANDOS PROPIOS
% =====================================================================

% --- Portada estándar ---------------------------------------------------
% Uso: \portada{\docTitulo}
\newcommand{\logoUniversidad}{img/logo-ucu.png}
\newcommand{\portada}[1]{%
  \begin{titlepage}
  \centering
  \vspace*{1.5cm}
  \IfFileExists{\logoUniversidad}{%
    \includegraphics[width=0.40\textwidth]{\logoUniversidad}\par
    \vspace{1.2cm}%
  }{%
    \fbox{\parbox[c][3cm][c]{0.5\textwidth}{\centering\color{leipAmbar}%
      \textbf{[Logo pendiente]}\\[4pt]\small\texttt{\logoUniversidad}}}\par
    \vspace{1.2cm}%
  }
  {\Large \proyectoUniversidad\par}
  {\large \proyectoFacultad\par}
  {\large \proyectoCarrera\par}
  \vspace{2.5cm}
  {\Huge\bfseries\color{leipAzul} #1\par}
  \vspace{0.8cm}
  {\Large \proyectoNombre{} --- \proyectoNombreLargo\par}
  \vspace{2cm}
  {\large Versi\'on \docVersion\par}
  {\large \docFecha\par}
  \vfill
  \begin{tabular}{c}
    Sebastian Forische \\
    Armando Hernandez de la Vega \\
    Luana Acu\~na \\
    Bruno Sebastian Olivera Viera Gomez \\
  \end{tabular}
  \vspace{1.5cm}
  \end{titlepage}
}

% --- Historial de revisiones -------------------------------------------
% Uso: \begin{historial} 1.0 & 26/06/2026 & Versión inicial & Armando \\ \end{historial}
% Nota: usa tabular con p{} y NO tabularx. tabularx necesita leer el cuerpo
% completo de la tabla de una sola vez, y no tolera que \begin y \end queden
% en macros distintas.
\newenvironment{historial}{%
  \section*{Historial de revisiones}
  \addcontentsline{toc}{section}{Historial de revisiones}
  \begin{tabular}{@{}l l p{0.52\textwidth} l@{}}
  \toprule
  \textbf{Versi\'on} & \textbf{Fecha} & \textbf{Descripci\'on} & \textbf{Autor} \\
  \midrule
}{%
  \bottomrule
  \end{tabular}
  \clearpage
}

% --- Marcador de contenido pendiente -----------------------------------
% Uso: \pendiente{resumen corto}{texto completo, puede traer listas}
%   - el resumen corto va al listado del final: debe ser texto plano
%   - el texto completo va al recuadro en el cuerpo del documento
% Para la entrega final: \pendientesfalse oculta los recuadros del PDF,
% pero el listado del final sigue mostrando qué quedó abierto.
\newif\ifpendientes
\pendientestrue

\makeatletter
\newcommand{\listofpendientes}{%
  \section*{Pendientes de este documento}
  \addcontentsline{toc}{section}{Pendientes de este documento}
  \@starttoc{pnd}%
}
\newcommand{\pendiente}[2]{%
  \addcontentsline{pnd}{section}{\protect\numberline{}#1}%
  \ifpendientes
    \begin{tcolorbox}[breakable,colback=leipAmbarBg,colframe=leipAmbar,
                      boxrule=0.6pt,left=6pt,right=6pt,top=4pt,bottom=4pt]
      \textbf{\color{leipAmbar}PENDIENTE.}\ #2
    \end{tcolorbox}
  \fi
}
\makeatother

% --- Identificadores de requerimientos ---------------------------------
% Se usan en requerimientos.tex; van acá para que la numeración sea única.
\newcounter{rfcounter}
\newcommand{\rf}[2]{% \rf{etiqueta}{texto}
  \refstepcounter{rfcounter}%
  \label{rf:#1}%
  \textbf{RF-\arabic{rfcounter}}\quad #2\par\smallskip
}
\newcounter{rnfcounter}
\newcommand{\rnf}[2]{%
  \refstepcounter{rnfcounter}%
  \label{rnf:#1}%
  \textbf{RNF-\arabic{rnfcounter}}\quad #2\par\smallskip
}

% --- Figura con fallback si la imagen todavía no existe ----------------
% Uso: \figuraLEIP{img/arquitectura.png}{Arquitectura de LEIP}{fig:arq}{0.9}
\newcommand{\figuraLEIP}[4]{%
  \begin{figure}[htbp]
  \centering
  \IfFileExists{#1}{\includegraphics[width=#4\textwidth]{#1}}{%
    \fbox{\parbox[c][4cm][c]{0.8\textwidth}{\centering\color{leipAmbar}%
      \textbf{[Figura pendiente]}\\[4pt]\small\texttt{#1}}}%
  }
  \caption{#2}
  \label{#3}
  \end{figure}
}

\begin{document}


\portada{\docTitulo}

\tableofcontents
\clearpage

\begin{historial}
0.1 & 08/09/2026 & Versi\'on inicial. Re\'une el contenido de mercado, validaci\'on comercial y modelo de negocio trasladado desde el Documento de Visi\'on & Equipo \\
0.2 & 12/09/2026 & Revisi\'on de redacci\'on: se elimina la escritura defensiva; el estado de la validaci\'on se enuncia una sola vez, en su propia secci\'on & Equipo \\
\end{historial}

% =====================================================================
\section{Introducci\'on}

\subsection{Prop\'osito}

Este documento re\'une los aspectos comerciales de \proyectoNombre{}
(\proyectoNombreLargo): el mercado al que se dirige, la validaci\'on del
problema desde el punto de vista comercial, el posicionamiento frente a las
alternativas existentes y el modelo de negocio con el que se sostiene.

La definici\'on del producto, sus capacidades, sus usuarios y sus restricciones
t\'ecnicas se desarrollan en el Documento de Visi\'on. Este documento parte de
esa definici\'on y no la repite. El contenido de partida proviene del Documento
de Visi\'on, de donde se traslad\'o sin modificaciones.

% =====================================================================
\section{Mercado}

\subsection{Mercado objetivo}

El mercado objetivo de \proyectoNombre{} est\'a compuesto por analistas, equipos
de investigaci\'on, tesorer\'ias, consultoras, fondos de inversi\'on, brokers,
acad\'emicos, investigadores y organizaciones que trabajan de forma recurrente
con informaci\'on econ\'omica y financiera de Am\'erica Latina.

Este mercado necesita reducir el tiempo dedicado a buscar, descargar, limpiar y
comparar informaci\'on de distintos organismos. Tambi\'en requiere seguir
revisiones hist\'oricas, mantener trazabilidad y relacionar expectativas,
inflaci\'on, actividad, riesgo fiscal, sector externo y condiciones financieras
dentro de un mismo entorno de an\'alisis.

\subsection{Cobertura geogr\'afica y prioridad comercial}

La cobertura inicial del producto ---Brasil, M\'exico y Uruguay--- responde
tambi\'en a un orden de prioridad comercial: Brasil es el mercado prioritario
por escala institucional, M\'exico el segundo mercado objetivo y Uruguay el
mercado de proximidad donde se realiza la validaci\'on inicial. Chile se
considera relevante, pero su incorporaci\'on queda bloqueada comercialmente
hasta obtener permiso de uso del organismo emisor.

% =====================================================================
\section{Validaci\'on comercial del problema}
\label{sec:validacion}

La prioridad inmediata es validar una pregunta central: si consolidar y comparar
expectativas macroecon\'omicas es una tarea recurrente y costosa para el
segmento que realmente paga, o si es una tarea espor\'adica que no justifica una
plataforma.

\subsection{Condiciones de problema v\'alido}

\begin{table}[htbp]
\centering
\caption{Condiciones de problema v\'alido}
\label{tab:validez}
\begin{tabularx}{\textwidth}{l X l}
\toprule
\textbf{Condici\'on} & \textbf{Aplicaci\'on en \proyectoNombre{}} & \textbf{Estado} \\
\midrule
Masivo dentro del nicho &
Debe afectar a suficientes analistas, equipos de investigaci\'on, consultoras y
tesorer\'ias que trabajen con expectativas macroecon\'omicas. &
Pendiente de muestra suficiente \\
\addlinespace
Doloroso &
Debe ser frecuente y costoso en tiempo como para justificar pago o adopci\'on
institucional. &
A validar con encuestas y entrevistas \\
\addlinespace
Mal resuelto &
Las soluciones actuales existen, pero son caras, gen\'ericas, manuales o no
est\'an centradas en expectativas oficiales. &
Hip\'otesis con evidencia secundaria \\
\bottomrule
\end{tabularx}
\end{table}

Ninguna de las tres condiciones est\'a validada: la muestra actual no alcanza.
La validaci\'on contin\'ua hasta reunir suficientes encuestas y entrevistas,
desagregadas por perfil institucional. Hasta entonces, el modelo de negocio de
la secci\'on \ref{sec:negocio} es una hip\'otesis.

\subsection{Plan de validaci\'on}

La validaci\'on se realiza mediante una encuesta dirigida a analistas
macroecon\'omicos, equipos de investigaci\'on, tesorer\'ias, consultoras,
fondos, acad\'emicos y perfiles que trabajen con expectativas econ\'omicas de
Am\'erica Latina. El objetivo es entender si el problema existe en la
pr\'actica: si consultar, comparar y seguir expectativas macro es una tarea
frecuente, lenta y mal resuelta.

La validaci\'on se enfoca en tres se\~nales:

\begin{enumerate}
  \item cu\'antas veces el usuario trabaj\'o con expectativas macro en los
        \'ultimos meses;
  \item cu\'anto tiempo le tom\'o conseguir y ordenar los datos;
  \item si necesita seguir revisiones en el tiempo o s\'olo consultar el dato
        puntual.
\end{enumerate}

Los resultados se analizan por tipo de usuario. Las respuestas de perfiles
institucionales tienen mayor peso, porque representan mejor al cliente que
podr\'ia pagar por el producto. Las respuestas de academia o usuarios
individuales sirven como aprendizaje, pero no son suficientes por s\'i solas
para validar el negocio.

El instrumento utilizado es la encuesta disponible en
\url{https://docs.google.com/forms/d/e/1FAIpQLSe4jNRspwc_cyIJii9npmMjpXSXsAxOBzT6vkGCEi5vW80DEQ/viewform}.

\subsection{Estado de la validaci\'on}

La recolecci\'on de respuestas est\'a en curso y la muestra disponible no
alcanza para validar el problema, la disposici\'on a pagar ni el modelo de
negocio. Esta versi\'on no presenta conclusiones cuantitativas.

Al cerrar la recolecci\'on, este documento reporta el n\'umero de respuestas, su
composici\'on por perfil y los tres indicadores del plan de validaci\'on. Las
conclusiones separan los perfiles institucionales de los acad\'emicos e
individuales, e incorporan la ronda de entrevistas cualitativas.

% =====================================================================
\section{Posicionamiento}

\subsection{Declaraci\'on de posicionamiento}

\begin{table}[htbp]
\centering
\begin{tabularx}{\textwidth}{>{\bfseries}l X}
\toprule
Para & analistas, equipos de investigaci\'on, tesorer\'ias, consultoras y fondos
que trabajan de forma recurrente con informaci\'on econ\'omica de Am\'erica Latina, \\
\addlinespace
Que & necesitan consultar, comparar y seguir expectativas y variables
macroecon\'omicas de varios pa\'ises sin mantener procesos propios de descarga,
limpieza y normalizaci\'on, \\
\addlinespace
\proyectoNombre{} es & una plataforma SaaS de inteligencia econ\'omica regional \\
\addlinespace
Que & transforma informaci\'on oficial dispersa en m\'etricas comparables y
trazables, con seguimiento de revisiones y cruces entre dominios econ\'omicos, \\
\addlinespace
A diferencia de & las terminales financieras generalistas, las bases macro
especializadas, los paneles privados en PDF y las planillas internas, \\
\addlinespace
Nuestro producto & se especializa en informaci\'on oficial de Am\'erica Latina y
en la din\'amica de sus revisiones, entregando comparabilidad regional lista
para usar a una fracci\'on del costo y la complejidad de las alternativas. \\
\bottomrule
\end{tabularx}
\caption{Declaraci\'on de posicionamiento}
\label{tab:posicionamiento}
\end{table}

\subsection{Alternativas y competencia}

El an\'alisis competitivo muestra que existen alternativas para obtener datos de
expectativas macroecon\'omicas, pero ninguna trabaja sobre expectativas
oficiales con foco en revisi\'on, divergencia y comparabilidad regional
exclusiva de Am\'erica Latina.

\begin{table}[htbp]
\centering
\caption{Productos similares y alternativas}
\label{tab:competencia}
\begin{tabularx}{\textwidth}{l l X}
\toprule
\textbf{Alternativa} & \textbf{Tipo} & \textbf{Por qu\'e no resuelve el problema} \\
\midrule
Fuentes oficiales directas & Manual &
Confiables, pero dispersas, heterog\'eneas y sin comparabilidad regional lista. \\
\addlinespace
Bloomberg, Refinitiv & Terminal financiera &
Amplias y robustas, pero extremadamente costosas y no enfocadas en la capa de
expectativas de Am\'erica Latina. \\
\addlinespace
CEIC, Macrobond, Haver & Bases macro especializadas &
Fuertes en infraestructura de investigaci\'on, pero costosas y con un foco
m\'as amplio. \\
\addlinespace
Consensus Economics & Panel privado de informes &
Producto distinto: panel privado en PDF, no un \emph{dataset} din\'amico y
explorable de expectativas. \\
\addlinespace
Trading Economics y similares & Datos macro observados &
Foco en datos observados, no en la din\'amica de revisiones de encuestas oficiales. \\
\addlinespace
Planillas internas y agentes de IA & Soluci\'on artesanal &
Resuelven parcialmente, pero requieren mantenimiento, control de calidad y
responsabilidad interna que consumen tiempo y esfuerzo. \\
\bottomrule
\end{tabularx}
\end{table}

\proyectoNombre{} se diferencia por transformar informaci\'on econ\'omica
oficial de Am\'erica Latina en m\'odulos comparables y conectados. No se limita
a visualizar datos: incorpora m\'etricas determin\'isticas, hist\'oricos,
alertas, cruces entre expectativas, inflaci\'on y actividad, y una lectura
recurrente sobre qu\'e se movi\'o en la regi\'on. Frente a herramientas
generalistas ofrece menor costo y complejidad; frente a planillas internas o
soluciones improvisadas con IA aporta trazabilidad, consistencia,
actualizaci\'on continua, seguridad e integraci\'on program\'atica.

% =====================================================================
\section{Modelo de negocio}
\label{sec:negocio}

El modelo combina acceso \emph{freemium}, contrataci\'on modular, planes
institucionales y servicios profesionales. La versi\'on gratuita permite
consultar las series oficiales que ya se encuentran disponibles p\'ublicamente,
con trazabilidad hacia su origen. Los planes pagos no cobran por el dato oficial
aislado, sino por el hist\'orico consolidado y versionado, las comparaciones, la
normalizaci\'on, las m\'etricas propias, los cruces, las alertas, los
\emph{briefings}, las exportaciones y el acceso program\'atico.

\subsection{Planes}

La tabla \ref{tab:planes} presenta los planes, su precio de referencia y el
valor que incluye cada uno.

\begin{table}[htbp]
\centering
\caption{Planes y precios de referencia}
\label{tab:planes}
\begin{tabularx}{\textwidth}{l l X}
\toprule
\textbf{Plan} & \textbf{Precio mensual} & \textbf{Valor incluido} \\
\midrule
Free              & Sin costo        & Radar resumido y acceso limitado a indicadores y series seleccionadas. \\
M\'odulo individual & USD 18 por m\'odulo & Acceso completo al m\'odulo elegido. Hasta tres m\'odulos contratables. \\
Completo          & USD 59           & Todos los m\'odulos, m\'etricas y cruces disponibles. \\
Team              & USD 349          & Cinco usuarios, plataforma completa, API institucional, administraci\'on centralizada y \emph{dashboard} configurable. \\
Enterprise        & Desde USD 799    & Soluci\'on configurable, entorno aislado, SLA, soporte prioritario e integraciones espec\'ificas. \\
\bottomrule
\end{tabularx}
\end{table}

Los m\'odulos individuales pueden contratarse hasta un m\'aximo de tres; cuando
el usuario necesita mayor cobertura, migra al plan Completo. Los cruces entre
dominios est\'an disponibles cuando la cuenta tiene acceso a todos los m\'odulos
involucrados. Los planes Team y Enterprise se orientan a organizaciones que
necesitan m\'ultiples usuarios, API institucional, administraci\'on
centralizada, aislamiento y soporte.

\subsection{Capa gratuita y capa paga}

La estrategia distingue el dato oficial de la elaboraci\'on propia. La capa
gratuita ofrece acceso consolidado a series que los organismos ya publican,
siempre con atribuci\'on y enlace a la fuente. La propuesta de valor comercial
se concentra en los hist\'oricos versionados, la normalizaci\'on
metodol\'ogica, las comparaciones regionales y las m\'etricas, se\~nales y
lecturas anal\'iticas desarrolladas por \proyectoNombre{}.

El radar gratuito cumple adem\'as una funci\'on comercial: permite a un usuario
evaluar la herramienta antes de comprometer presupuesto, y act\'ua como capa de
entrada hacia los planes pagos.

\subsection{Costos y viabilidad}

La estimaci\'on supone una primera etapa con hasta 500 usuarios, ocho actualizaciones mensuales, 50~GB de almacenamiento, 12 horas
mensuales de mantenimiento t\'ecnico, 10 horas mensuales de administraci\'on y
un valor estimado del trabajo de USD 15 por hora.

\begin{table}[htbp]
\centering
\caption{Costos operativos estimados}
\label{tab:costos}
\begin{tabular}{l r r}
\toprule
\textbf{Concepto} & \textbf{Mensual} & \textbf{Anual} \\
\midrule
Infraestructura           & USD 90  & USD 1.080 \\
LLM                       & Por estimar & Por estimar \\
Mantenimiento t\'ecnico   & USD 180 & USD 2.160 \\
Administraci\'on operativa & USD 150 & USD 1.800 \\
\midrule
\textbf{Subtotal sin LLM} & \textbf{USD 420} & \textbf{USD 5.040} \\
\bottomrule
\end{tabular}
\end{table}

La viabilidad b\'asica se eval\'ua con usuarios individuales y planes completos,
sin depender de vender un contrato Enterprise.

\pendiente{Estimar el costo del LLM a partir de la arquitectura y el volumen real}{Calcular el costo cuando est\'en definidos el primer boceto de
arquitectura de agentes, los modelos candidatos y el inventario inicial de
fuentes. La estimaci\'on debe considerar: ejecuciones por fuente y frecuencia,
tokens de entrada y salida, llamadas por extracci\'on, reintentos, validaciones,
generaci\'on de briefings y escenarios base, conservador y de estr\'es.}

\pendiente{Cerrar la validaci\'on comercial y actualizar el modelo de negocio}{Completar la recolecci\'on de encuestas y la ronda de entrevistas,
reportar los tres indicadores del plan de validaci\'on desagregados por perfil,
y revisar planes, precios y segmentos a partir de esa evidencia.}

% =====================================================================
\clearpage
\listofpendientes

\end{document}
